\documentclass[12pt,a4paper]{article}
\usepackage[utf8]{inputenc}
\usepackage[T2A]{fontenc}
\usepackage[russian]{babel}
\usepackage{amsmath,amssymb}
\usepackage{array}
\usepackage{longtable}
\usepackage{booktabs}
\usepackage{geometry}
\geometry{margin=2.3cm}

\title{Протокол исследования по сравнению корпусов \texttt{TOE.pdf} и \texttt{ugsm\_dynamics\_audit-3.pdf}}
\author{}
\date{}

\begin{document}

\maketitle

\begin{abstract}
Настоящий протокол фиксирует, что именно было проверено в рамках сравнительного исследования двух теоретических корпусов --- \texttt{TOE.pdf} и \texttt{ugsm\_dynamics\_audit-3.pdf}, какие структурные и методологические результаты были найдены, а также к каким итоговым выводам приводит проведённый анализ. Документ написан в формате исследовательского отчёта: цель, объекты проверки, карта проведённых проверок, основные находки, итоговые выводы и практические рекомендации для дальнейшего усиления программы.
\end{abstract}

\section{Цель исследования}

Целью исследования было не просто сопоставить две теоретические работы по содержанию, а определить:
\begin{itemize}
    \item насколько каждая из них продвинута как исследовательская программа;
    \item где находятся их реальные точки пересечения;
    \item где именно проходят различия по строгости, воспроизводимости и зрелости;
    \item каков статус открытых проблем в каждом корпусе;
    \item какая из программ на текущем этапе лучше подготовлена к режиму S2T-совместимого аудита.
\end{itemize}

Иными словами, исследование было направлено не только на сравнение идей, но и на сравнение качества научной организации этих идей.

\section{Объекты исследования}

В исследовании рассматривались два корпуса:
\begin{enumerate}
    \item \texttt{TOE.pdf} вместе с развивающим его архивом тематических документов;
    \item \texttt{ugsm\_dynamics\_audit-3.pdf} как аудит-документ, внутри которого теория UGSM доводится от эвристического уровня к более строгому динамическому и проверяемому статусу.
\end{enumerate}

Важно, что оба корпуса анализировались не как изолированные статичные тексты, а как развивающиеся исследовательские программы. Для TOE это особенно важно из-за распределения усилений по нескольким архивным документам. Для UGSM это важно потому, что внутри самого audit-документа ранние gap-блоки по ходу текста частично закрываются или переводятся в более сильный операциональный статус.

\section{Что именно было проверено}

В ходе исследования были проверены следующие блоки.

\subsection{Проверка концептуального профиля}

Было проверено, каков основной тип каждой программы:
\begin{itemize}
    \item является ли работа фундаментальной объединяющей архитектурой;
    \item является ли она локальным, но более дисциплинированным динамическим каркасом;
    \item каков её масштаб объяснения и степень охвата физических секторов.
\end{itemize}

\subsection{Проверка доказательного статуса}

Отдельно проверялось:
\begin{itemize}
    \item где текст говорит языком строгого вывода;
    \item где реально присутствует только частичная derivation-схема;
    \item как сами документы маркируют незакрытые узлы;
    \item различают ли они операциональную валидацию и финальное теоретическое замыкание.
\end{itemize}

\subsection{Проверка процедурной зрелости}

Для обеих программ оценивалось наличие или отсутствие следующих признаков:
\begin{itemize}
    \item train/test разделения;
    \item blind-валидации;
    \item negative-control логики;
    \item baseline-сравнений;
    \item фиксированных критериев прохождения цикла;
    \item встроенной anti-overfit дисциплины.
\end{itemize}

\subsection{Проверка прогрессивной структуры корпуса}

Было проверено, развивается ли каждая теория поступательно:
\begin{itemize}
    \item закрываются ли ранние пробелы в поздних разделах или поздних документах;
    \item локализованы ли открытые проблемы явно;
    \item существует ли внутренняя карта усиления теории.
\end{itemize}

\subsection{Проверка текущей научной зрелости}

Для итогового сравнения проверялись:
\begin{itemize}
    \item воспроизводимость постановки;
    \item локальная математическая дисциплина;
    \item глобальная замкнутость;
    \item честность самокритики;
    \item фальсифицируемость;
    \item готовность к внешнему и внутреннему аудиту.
\end{itemize}

\section{Что было найдено}

\subsection{Главная находка по TOE}

По корпусу \texttt{TOE.pdf} было найдено следующее:
\begin{itemize}
    \item это программа очень высокой архитектурной амбиции, претендующая на широкую унификацию фундаментальной физики;
    \item её сильнейшая сторона --- масштаб замысла, число охватываемых секторов и потенциальная фундаментальная значимость;
    \item её ключевая слабость --- неполная процедурная стандартизация как единого audit-контура;
    \item часть слишком общих ранних тезисов в архиве действительно получает более локальные усиления, proof-fragments и roadmaps;
    \item при этом корпус остаётся неоднородным: усиления распределены по нескольким документам, а не собраны в одном жёстко организованном протоколе проверки.
\end{itemize}

\subsection{Главная находка по UGSM}

По корпусу \texttt{ugsm\_dynamics\_audit-3.pdf} было найдено следующее:
\begin{itemize}
    \item программа существенно уже по охвату, чем TOE, но заметно сильнее по процедурной дисциплине;
    \item документ специально организован как audit-доведение от эвристики к минимальному динамическому EFT-каркасу;
    \item в нём явно встроены blind train/test-подход, negative controls и критерии прохождения;
    \item документ честно различает статус ``operationally validated'' и статус полностью завершённого физического вывода;
    \item внутри текста прослеживается прогрессия: часть ранних gap-блоков действительно сужается, локально закрывается или переводится в более строгую форму.
\end{itemize}

\subsection{Общая находка по пересечениям}

Исследование показало, что между двумя корпусами есть реальные глубокие пересечения:
\begin{enumerate}
    \item обе программы опираются на спектрально-геометрический тип мышления;
    \item обе стремятся связать наблюдаемые с операторной или спектральной структурой;
    \item обе признают, что численные совпадения сами по себе недостаточны;
    \item обе движутся от эвристики к более строгой постановке;
    \item обе в той или иной степени приближаются к S2T-совместимому режиму оценки.
\end{enumerate}

\subsection{Главная находка по различиям}

Наиболее важное различие оказалось не только содержательным, но и организационным:
\begin{itemize}
    \item TOE сильнее как большая объединяющая архитектура;
    \item UGSM сильнее как процедурно зрелый локальный аудит;
    \item TOE выигрывает по горизонту и предельной амбиции;
    \item UGSM выигрывает по прозрачности, воспроизводимости и anti-overfit дисциплине на текущем этапе.
\end{itemize}

\section{Таблица исследовательских находок}

\begin{longtable}{p{5.2cm}p{4.1cm}p{4.1cm}}
\toprule
Что проверялось & Что найдено в TOE & Что найдено в UGSM \\
\midrule
\endhead
Масштаб программы & Почти полная унификационная архитектура & Узкий, но управляемый EFT-аудит \\

Ясность статуса утверждений & Не всегда равномерна в базовых текстах; усиливается в архиве & Очень явно маркируется внутри основного документа \\

Процедурная дисциплина & Частично присутствует, но не собрана в единый стандартный контур & Сильная train/test/blind/negative-control организация \\

Структура развития & Идёт через распределённый архив тематических усилений & Идёт внутри одного audit-документа как последовательное доведение \\

Открытые проблемы & Признаны и частично локализованы через roadmap-документы & Признаны явно, стратифицированы и связаны с финальным DoD \\

Воспроизводимость & Средняя, зависит от конкретного блока архива & Высокая на уровне audit-процедуры \\

Готовность к S2T-аудиту & Растущая, но неоднородная & Очень высокая по текущему формату документа \\
\bottomrule
\end{longtable}

\section{Основные выводы исследования}

По итогам проведённой проверки были сделаны следующие выводы.

\begin{enumerate}
    \item \texttt{TOE.pdf} и \texttt{ugsm\_dynamics\_audit-3.pdf} нельзя корректно сравнивать только по масштабу идей; их необходимо сравнивать и по уровню процедурной зрелости.
    \item TOE на данный момент выглядит сильнее как фундаментальная объединяющая программа, но слабее как единый стандартизированный audit-контур.
    \item UGSM на данный момент выглядит сильнее как исследовательский документ: более воспроизводимый, более дисциплинированный и более готовый к честной верификации.
    \item Сильный результат UGSM не означает, что она фундаментальнее TOE; это означает, что она лучше доведена до формата контролируемого научного исследования здесь и сейчас.
    \item Более слабая процедурная позиция TOE не означает слабость самой идеи; она означает, что масштаб программы пока опережает степень её формальной и audit-организации.
    \item В обоих случаях была обнаружена важная прогрессивная динамика: и TOE, и UGSM не являются статичными текстами, а представляют собой программы постепенного доведения и локального закрытия промежуточных узлов.
\end{enumerate}

\section{Итоговый вердикт}

Итоговый вердикт исследования можно сформулировать так:

\begin{quote}
На текущем этапе UGSM выглядит более зрелой как исследовательская машина и как объект методологической проверки, тогда как TOE остаётся более мощной по горизонту, масштабу и потенциальному фундаментальному значению в случае успешного замыкания.
\end{quote}

Иными словами, если критерий сравнения --- текущая научная дисциплина, воспроизводимость и S2T-совместимость, преимущество находится на стороне UGSM. Если же критерий сравнения --- предельная амбиция и возможный масштаб теоретического прорыва, преимущество остаётся за TOE.

\section{Практический вывод по итогам исследования}

Исследование показывает, что для дальнейшего усиления корпуса TOE наиболее полезны следующие шаги:
\begin{itemize}
    \item внедрение более явного train/test разделения там, где это возможно;
    \item введение blind- и negative-control режимов для численных и феноменологических блоков;
    \item единая маркировка статусов вида ``operational'', ``derived'', ``open'', ``roadmap'';
    \item перенос audit-логики из периферийных материалов в основной корпус;
    \item сборка единого документа, в котором открытые узлы и степень их закрытия были бы представлены столь же прозрачно, как в UGSM.
\end{itemize}

Для UGSM главный путь усиления противоположен: расширить локально успешный audit-контур до более широкого архитектурного охвата и связать частные спектральные достижения с более цельной фундаментальной картиной.

\section{Заключение}

В рамках данного исследования было проверено не только содержание двух корпусов, но и их форма как научных программ. Было найдено, что TOE и UGSM движутся по частично общему спектрально-операторному полю, но реализуют две разные стратегии теоретического развития. TOE --- это стратегия максимальной объединяющей архитектуры с распределённым архивом усилений. UGSM --- это стратегия локального, но дисциплинированного audit-доведения с высокой процедурной прозрачностью. Главный вывод состоит в том, что на текущем этапе UGSM сильнее по научной зрелости исследовательского формата, тогда как TOE сохраняет преимущество по масштабу замысла и потенциальной фундаментальной значимости. Именно это различие и должно быть положено в основу дальнейшего развития обеих программ.

\end{document}